\documentclass[aps,pra,preprint,superscriptaddress]{revtex4-2}
\usepackage{graphicx}
\usepackage{amsmath,amssymb,bm}
\usepackage{hyperref}
\hypersetup{colorlinks=true,linkcolor=blue,citecolor=blue,urlcolor=blue}

\begin{document}

\title{Measurement-Induced Criticality and Non-Markovian Memory in an
Event-Driven Quantum-Trajectory Simulator}

\author{The RCT Collaboration}
\affiliation{Open-science telemetry release.
Correspondence: \url{https://github.com/hackall360/rct-quantum-lattice-results}}
\date{November 7, 2026}

\begin{abstract}
We report two hardware-backed experiments performed with an event-driven
quantum-trajectory simulator on a single consumer GPU: (i) a
measurement-induced phase transition in a 16-qubit brickwork circuit, where
the mean half-cut entropy falls from a volume-law plateau of $7.61$ bits
(95\% of maximum) to an area-law floor of $2.31$ bits with a mid-entropy
crossing at $p_c \approx 0.150$ and peak fluctuations at the crossing; and
(ii) five Breuer--Piilo--Lorenzo-certified information-backflow revivals
(largest $+0.198$) in the trace distance of two core preparations sharing a
structured bath.  Both witnesses are computed per trajectory with ensemble
averaging last, an ordering that mixed-state averaging would destroy.  All
numbers are reproduced from two CSV files by a single dependency-free
plotting script; no proprietary code is needed to verify any claim.
\end{abstract}

\maketitle

\section{Introduction}

The competition between entangling unitary dynamics and disentangling
measurement produces a purely dynamical phase transition in the structure of
individual quantum trajectories~\cite{LiChenFisher2018,LiChenFisher2019,
SkinnerRuhmanNahum2019,ChanNandkishore2019,Choi2020,FisherReview2023}: below
a critical measurement rate $p_c$ the steady state obeys a volume law for the
entanglement entropy, above it an area law.  Observing the transition requires
statistics of \emph{pure trajectories} --- averaging the states first into a
density matrix erases the very ordering the transition is defined
by~\cite{FisherReview2023}.  The same ordering underlies witnesses of
non-Markovianity: the trace-distance contractivity theorem of
Breuer, Piilo, and Lorenzo certifies information backflow from environment to
system whenever the distinguishability of two preparations
rises~\cite{BLP2009,RivasHuelgaPlenio2014,deVegaAlonso2017}, a signal invisible
to any Markovian description and expensive to capture with exact hierarchical
methods~\cite{TanimuraKubo1989,Strathearn2018}.

This paper documents two such observations obtained from one simulation
engine: a quantum-trajectory Monte Carlo program~\cite{Carmichael1993,
DalibardCastinMolmer1992,PlenioKnight1998} in which wavepacket histories
propagate through a three-dimensional geometry, interact at coincident-arrival
events, and collapse under Born-rule branching.  Per-trajectory observables
are evaluated first; ensemble averaging is performed last.  We report the
empirical record at the macro level --- inputs, outputs, and the script that
connects them --- and deliberately exclude implementation internals.

\section{Method in brief}

\paragraph{Trajectory ensemble.}
Each trajectory is a wavepacket history propagated through a culled 3D
geometry.  Coincident arrivals interact through two-body unitaries; a Monte
Carlo branching step then selects survival versus collapse with Born-rule
statistics.  A trajectory record carries the event location, an observable in
$[-1,1]$, and a collapse type.

\paragraph{Ordering of averages.}
For every witness $W$, we compute $W$ on each pure trajectory and average the
scalars: $\bar W = \langle W(|\psi_k\rangle) \rangle_k$.  No state averaging
enters either protocol.

\paragraph{Hardware.}
All runs reported here executed on a single NVIDIA RTX~4070~SUPER consumer
GPU.  The release archive contains only the resulting numbers and the script
that plots them; verification requires Python~3 and matplotlib, no GPU and no
proprietary dependency.

\section{Experiment 1: 3D measurement-induced phase transition}

The circuit acts on $N=16$ qubits with 48 layers, each layer applying 8
Haar-random two-qubit gates~\cite{Mezzadri2007} on disjoint pairs, with
measurement placements drawn from the culled 3D geometry at per-qubit
per-layer rates $p \in [0.02,0.40]$, 8 shots per point.  The witness is the
mean half-cut entanglement entropy $\bar S$ in bits (maximum 8).

\begin{figure}[t]
\includegraphics[width=\columnwidth]{fig_mipt.png}
\caption{Mean half-cut entropy versus measurement rate ($N=16$, 48 layers
$\times$ 8 Haar-random gates, 8 shots per point).  Dashed lines mark the
8-bit maximum and the mid-entropy crossing $p_c \approx 0.150$.  Error bars
are one sample standard deviation; note the fluctuation peak at the
crossing. \label{fig:mipt}}
\end{figure}

\begin{table}[t]
\caption{MIPT sweep: measurement rate $p$, mean half-cut entropy $\bar S$,
sample standard deviation $\sigma$, and mean measurements per trajectory.
\label{tab:mipt}}
\begin{tabular}{ccccc}
\hline\hline
$p$ & $\bar S$ (bits) & $\sigma$ & meas./traj & regime \\
\hline
0.02 & 7.609 & 0.354 & 13 & volume-law plateau \\
0.05 & 6.953 & 0.693 & 40 & plateau edge \\
0.08 & 7.022 & 0.451 & 58 & plateau edge \\
0.12 & 5.926 & 0.422 & 89 & descending \\
0.16 & 4.644 & 1.437 & 121 & crossing, peak fluctuations \\
0.20 & 4.282 & 1.094 & 156 & descending \\
0.25 & 3.776 & 0.621 & 191 & descending \\
0.30 & 2.458 & 0.977 & 230 & floor onset \\
0.35 & 2.397 & 0.820 & 268 & floor \\
0.40 & 2.310 & 0.528 & 303 & area-law floor \\
\hline\hline
\end{tabular}
\end{table}

The entropy descends from $7.609 \pm 0.354$ bits (95\% of the 8-bit maximum)
to $2.310 \pm 0.528$ bits.  The mid-entropy crossing sits at
$p_c \approx 0.150$, and the sample variance peaks at the crossing
($\pm 1.44$ at $p=0.16$), the finite-size signature of critical fluctuations.
A random-unitary oracle cross-check of the entropy estimator agrees to
$9\times 10^{-3}$ bits.

\paragraph{Limits stated plainly.}
$N=16$ locates a crossover, not the thermodynamic limit; no finite-size
scaling collapse is claimed.  Each point uses 8 shots.  Measurement
accounting is per gate, not per wall-clock hit.  These caveats bound the
claim to what the data support: a finite-size MIPT signature with the
expected plateau, drop, floor, and fluctuation peak.

\section{Experiment 2: non-Markovian backflow}

A 3-qubit core is entangled with a 6-site structured bath and evolved under
a fixed event script.  Two core preparations are tracked; their trace
distance $D(t)$ is compared against contractivity.  Any rise in $D$ is a
theorem-backed certificate of non-Markovianity~\cite{BLP2009} that no
Markovian channel can forge.  A forgetful (repeatedly reset) arm is shown as
an informational control; because the reset channels are state-dependent,
cross-arm distances carry no contractivity obligation, and the rigorous
witness is the memory arm alone.

\begin{figure}[t]
\includegraphics[width=\columnwidth]{fig_blp.png}
\caption{Trace distance $D(t)$ of two core preparations versus event step.
Solid: memory arm --- five rises, largest $+0.198$ ($0.664 \to 0.862$).
Dashed: forgetful control (informational only). \label{fig:blp}}
\end{figure}

\begin{table}[t]
\caption{BLP series: event step, memory-arm $D$, forgetful-control $D$.
Rises of the memory arm occur at steps $4, 7, 8, 11, 13$.
\label{tab:blp}}
\begin{tabular}{ccc}
\hline\hline
step & $D_{\mathrm{memory}}$ & $D_{\mathrm{control}}$ \\
\hline
0 & 1.000 & 1.000 \\
1 & 1.000 & 1.000 \\
2 & 0.959 & 0.906 \\
3 & 0.950 & 0.881 \\
4 & 0.959 & 0.889 \\
5 & 0.909 & 0.879 \\
6 & 0.834 & 0.835 \\
7 & 0.871 & 0.922 \\
8 & 0.900 & 0.851 \\
9 & 0.797 & 0.818 \\
10 & 0.664 & 0.816 \\
11 & 0.862 & 0.924 \\
12 & 0.838 & 0.773 \\
13 & 0.934 & 0.798 \\
\hline\hline
\end{tabular}
\end{table}

The memory arm exhibits \textbf{five distinct rises} with maximum
single-step revival \textbf{$+0.198$} ($0.664 \to 0.862$ at steps
$10 \to 11$), including a late-evolution recovery that rules out transient
artifacts.  The forgetful control shows no comparable coherent revival.  The
reset channel underlying the control arm is validated independently to
$10^{-16}$ plus idempotence.  A scale check with the same engine sustains the
bath at 20{,}000 modes with zero truncation error and exact norm
conservation.

\section{Discussion}

The two results share one methodological root: per-trajectory evaluation
followed by ensemble averaging.  A density-matrix-averaged ensemble would
report a single mixed state in both protocols and neither witness would
survive.  The MIPT curve and the BLP revivals are therefore joint evidence
that trajectory-resolved statistics reach regimes --- 3D criticality and
deep-memory transport --- that state-averaged methods cannot see, consistent
with the established understanding of trajectory phases~\cite{FisherReview2023}
and non-Markovian witnesses~\cite{RivasHuelgaPlenio2014}.

What is \emph{not} claimed: thermodynamic critical exponents, a scaling
collapse, or a general-purpose non-Markovian solver.  What \emph{is} claimed:
a finite-size 3D MIPT signature and five BLP-certified backflow revivals,
each number reproducible from the archived CSVs.

\section{Reproducibility}

The release archive provides \texttt{data/mipt\_sweep.csv} (Table~\ref{tab:mipt}),
\texttt{data/blp\_series.csv} (Table~\ref{tab:blp}),
\texttt{scripts/plot\_telemetry.py} (pure matplotlib, Agg backend, no network),
and the two figures.  Running
\begin{verbatim}
python3 scripts/plot_telemetry.py
\end{verbatim}
regenerates Fig.~\ref{fig:mipt} and Fig.~\ref{fig:blp} exactly.

\begin{acknowledgments}
Compute was performed on a single self-funded consumer GPU.  The authors thank
the developers of the open numerical and typesetting stack (Python,
matplotlib, TeX~Live, REVTeX) that makes this release verifiable.
\end{acknowledgments}

\begin{thebibliography}{99}
\bibitem{LiChenFisher2018} Y.~Li, X.~Chen, and M.~P.~A.~Fisher,
Phys.\ Rev.\ B \textbf{98}, 205136 (2018).
\bibitem{LiChenFisher2019} Y.~Li, X.~Chen, and M.~P.~A.~Fisher,
Phys.\ Rev.\ B \textbf{100}, 134306 (2019).
\bibitem{SkinnerRuhmanNahum2019} B.~Skinner, J.~Ruhman, and A.~Nahum,
Phys.\ Rev.\ X \textbf{9}, 031009 (2019).
\bibitem{ChanNandkishore2019} A.~Chan, R.~M.~Nandkishore, M.~Pretko, and
G.~Smith, Phys.\ Rev.\ B \textbf{99}, 224307 (2019).
\bibitem{Choi2020} S.~Choi, Y.~Bao, X.-L.~Qi, and E.~Altman,
Phys.\ Rev.\ Lett.\ \textbf{124}, 170502 (2020).
\bibitem{FisherReview2023} M.~P.~A.~Fisher, V.~Khemani, A.~Nahum, and
S.~Vijay, Annu.\ Rev.\ Condens.\ Matter Phys.\ \textbf{14}, 335 (2023).
\bibitem{BLP2009} H.-P.~Breuer, E.-M.~Laine, and J.~Piilo,
Phys.\ Rev.\ Lett.\ \textbf{103}, 210401 (2009).
\bibitem{RivasHuelgaPlenio2014} \'{A}.~Rivas, S.~F.~Huelga, and
M.~B.~Plenio, Rep.\ Prog.\ Phys.\ \textbf{77}, 094001 (2014).
\bibitem{deVegaAlonso2017} I.~de Vega and D.~Alonso,
Rev.\ Mod.\ Phys.\ \textbf{89}, 015001 (2017).
\bibitem{TanimuraKubo1989} Y.~Tanimura and R.~Kubo,
J.\ Phys.\ Soc.\ Jpn.\ \textbf{58}, 101 (1989).
\bibitem{Strathearn2018} A.~Strathearn, P.~Kirton, D.~Kilda, J.~Keeling, and
B.~W.~Lovett, Nat.\ Commun.\ \textbf{9}, 3322 (2018).
\bibitem{Carmichael1993} H.~Carmichael, \textit{An Open Systems Approach to
Quantum Optics} (Springer, Berlin, 1993).
\bibitem{DalibardCastinMolmer1992} J.~Dalibard, Y.~Castin, and K.~M\o lmer,
Phys.\ Rev.\ Lett.\ \textbf{68}, 580 (1992).
\bibitem{PlenioKnight1998} M.~B.~Plenio and P.~L.~Knight,
Rev.\ Mod.\ Phys.\ \textbf{70}, 101 (1998).
\bibitem{Mezzadri2007} F.~Mezzadri, Notices Amer.\ Math.\ Soc.\
\textbf{54}, 592 (2007).
\end{thebibliography}

\appendix
\section{Archived numbers}

For audit convenience the plotted numbers are the CSV contents verbatim.
MIPT $(p, \bar S, \sigma, \mathrm{meas/traj})$:
(0.02, 7.609, 0.354, 13), (0.05, 6.953, 0.693, 40),
(0.08, 7.022, 0.451, 58), (0.12, 5.926, 0.422, 89),
(0.16, 4.644, 1.437, 121), (0.20, 4.282, 1.094, 156),
(0.25, 3.776, 0.621, 191), (0.30, 2.458, 0.977, 230),
(0.35, 2.397, 0.820, 268), (0.40, 2.310, 0.528, 303).
BLP memory-arm $D(t)$, steps 0--13:
1.000, 1.000, 0.959, 0.950, 0.959, 0.909, 0.834, 0.871, 0.900, 0.797,
0.664, 0.862, 0.838, 0.934.

\end{document}
